problem:  
Let \( M^n \) be a compact, simply connected Riemannian manifold with **nonnegative sectional curvature** carrying an effective, isometric **torus action** \(T^k \curvearrowright M\).  Suppose \(k = \lfloor n/2 \rfloor - 1\), i.e. the torus rank is just *below* maximal symmetry rank, and assume that the action has **nonempty fixed‐point set**.  Let \(Q = M/T^k\) be the orbit space, equipped with its canonical Alexandrov metric of curvature \(\ge 0\).

Known examples with this symmetry range (such as products of spheres with diagonal \(T^k\)‐actions and biquotients) exhibit orbit spaces that are **manifolds with corners**, each codimension‑one face corresponding to a component of the fixed‐point set of a circle subgroup of \(T^k\); these faces are always **homeomorphic to disks**.

> **Precise open problem (Near‑maximal Symmetry Orbit Space Rigidity Conjecture):**  
> For \(M^n\) as above with \(k = \lfloor n/2 \rfloor - 1\),  
> - Must the orbit space \(Q = M/T^k\) be a **topological manifold with corners** whose faces are **all contractible**?  
> - Equivalently, can the boundary stratification of \(Q\) ever contain a face with nontrivial topology under nonnegative curvature?  
>  
> More generally, does nonnegative curvature enforce the following homological rigidity:  
> \[
> \widetilde{H}_i(\operatorname{Face}(Q)) = 0 \quad \text{for all codimension‑\(1\) faces and } 0 < i < \dim(\operatorname{Face}(Q)) ?
> \]  
>  
> A counterexample would be a nonnegatively curved \(M^n\) with a torus action just below maximal rank whose orbit space \(Q\) contains a topologically nontrivial face—e.g., a face homeomorphic to \(S^p \times D^q\) or with nontrivial fundamental group—something not yet known to exist.

Why is it a "good" problem:  

1. **Mathematical focus and plausibility:**  
   This version fixes the symmetry rank and the type of topology under study.  It is neither trivially false nor settled by known theorems.  The interplay between curvature and the topology of orbit faces is unexplored when one drops a single torus dimension from the maximal case.  The conjecture asks whether nonnegative curvature keeps the same rigourous control on orbit‑space topology or whether new behaviors arise.

2. **Bridging geometric and topological invariants:**  
   The problem connects **sectional curvature constraints** (a smooth, metric property) with the **combinatorial stratification** of the orbit space produced by a torus action.  Resolving it will require ideas from *Alexandrov geometry* (spaces of curvature ≥ 0), *transformation group theory*, and *topological data of group actions*.  It stands at the junction of these fields.

3. **Simple formulation, deep implications:**  
   The short question—“Can an orbit face of a nonnegatively curved torus action just below maximal rank be topologically nontrivial?”—is easy to state but touches the heart of the curvature–symmetry–topology interplay.  An affirmative or negative answer would reveal how close the near‑maximal case is to the rigid structure of maximal symmetry spaces.

4. **Extension beyond known results:**  
   Grove–Searle‑type results handle maximal rank.  Little is known one step below that bound.  The problem therefore precisely targets the first unclassified level in the hierarchy of torus actions under curvature bounds.

5. **Potential for new rigidity principles:**  
   A positive answer would extend curvature‑controlled “face contractibility” one symmetry level below maximal, hinting at a new quantitative rigidity principle linking curvature and torus action rank.  A counterexample could uncover a genuinely new orbit‑space topology, introducing new construction methods for nonnegatively curved manifolds.

Thus this narrowly focused conjecture—asking whether **nonnegative curvature enforces contractible faces in orbit spaces for torus actions just below maximal rank**—constitutes a concrete, open, and deep problem aligning with the essential frontier of curvature and symmetry interactions in Riemannian geometry.